\section{Evaluated Approaches}
\label{sec:approaches}

\paragraph{Feature engineering.}
We started from the state at turn 30 taken at face value: Pok\'emon alive and their HP, statuses,
the active Pok\'emon, boosts and Spikes. This scored 0.675. Adding Player~1's species and items and
the difference of move counts between the players gave a baseline of 0.685.

This state still counts Player~2's hidden knockouts as alive. How can we find them? Player~1's
knockouts are always visible, so every Player~1 Pok\'emon that leaves and never comes back (a
``vanishing'') is a real switch. If both players switch alike, the extra vanishings of Player~2 are
the hidden knockouts. A LightGBM classifier separates Player~2's vanishings from Player~1's using
only the event (HP before leaving, whether the replacement entered unhurt, move powers). With $q$ its output, the share of real switches among Player~2's
events is $(1-q)/q$, so
\[
P(\mathrm{KO}\mid x)=1-c\,\frac{1-q}{q},
\]
where $c$ makes the total match the observed excess. No labels are used, and Player~2 now loses as
many Pok\'emon as Player~1 (1.61 vs.\ 1.69 per battle). We recompute the state weighting each
Pok\'emon by its probability of being alive.

What about Player~1's identity? We add the smoothed win rate of Player~1's exact team, as a target
encoding computed out-of-fold inside each training fold.

\paragraph{Models.}
Our learner is LightGBM with shallow trees; logistic regression was 0.5 points worse, and XGBoost,
CatBoost and blends stayed within $\pm0.003$.

\paragraph{Validation.}
With differences this small, how do we know a block really helps? Every configuration uses the same
stratified 5-fold CV repeated three times, and we compare paired fold differences with the
Nadeau--Bengio corrected standard error~\cite{nadeau2003}, keeping a block only if it helps. This
dropped Player~2's estimated bench, extra state features and a damage formula.
